\chapter{Regularização e Seleção de Variáveis}
\label{cap:regularizacao}
% ── índice ──────────────────────────────────────────────────────
\index{regularização|textbf}
\index{LASSO|textbf}
\index{lasso@\texttt{lasso()}|textbf}
\index{Ridge|textbf}
\index{ridge_reg@\texttt{ridge\_reg()}|textbf}
\index{Elastic Net|textbf}
\index{enet@\texttt{enet()}|textbf}
\index{esparsidade|textbf}
\index{penalização L1 (lasso)}
\index{penalização L2 (ridge)}
\index{seleção de variáveis}

% ─────────────────────────────────────────────────────────────────────────────
\section{O problema de dimensionalidade}

Com $k$ regressores e $N$ pequeno, o OLS sofre de \emph{overfitting}: ajusta
o ruído da amostra além do sinal, elevando o erro fora da amostra. Quando
$k > N$, o OLS nem sequer é identificado. A regularização penaliza a
magnitude dos coeficientes, trocando viés por variância reduzida.

\subsection*{LASSO}

O LASSO (Tibshirani 1996) minimiza:

\[
  \hat\beta^{\mathrm{LASSO}} = \arg\min_\beta \left\{ \sum(y_i - x_i'\beta)^2
  + \lambda \sum_{j=1}^k |\beta_j| \right\}
\]

A penalidade $\ell_1$ força coeficientes pequenos a zero exatamente, fazendo
\emph{seleção de variáveis} automática. O caminho de regularização — LASSO
para diferentes $\lambda$ — revela quais variáveis entram conforme $\lambda$
decresce de $\infty$ (modelo nulo) a 0 (OLS).

\subsection*{Ridge}

A penalidade $\ell_2$ do Ridge (Hoerl e Kennard 1970) encolhe coeficientes
em direção a zero sem zerá-los:

\[
  \hat\beta^{\mathrm{Ridge}} = (X'X + \lambda I)^{-1} X'y
\]

Útil quando muitas variáveis têm pequeno efeito (contínuo de preditores
fracos). Não realiza seleção de variáveis.

\subsection*{Elastic Net}

O Elastic Net (Zou e Hastie 2005) combina $\ell_1$ e $\ell_2$:

\[
  \min_\beta \left\{ \sum(y_i - x_i'\beta)^2
  + \lambda \left[\rho\sum|\beta_j| + (1-\rho)\sum\beta_j^2\right] \right\}
\]

Herda seleção de variáveis do LASSO e estabilidade em grupos correlacionados
do Ridge.

% ─────────────────────────────────────────────────────────────────────────────
\section{Verificação por DGP}

DGP com 10 regressores, apenas $x_1$ e $x_2$ com efeito não-nulo.
$N = 200$ — regime onde regularização importa.

\begin{lstlisting}[caption={DGP esparso — 2 de 10 variáveis ativas}]
seed(42)
// 10 regressores; y = 1 + 2*x1 + 3*x2 + 0*(x3..x10) + e
generate df x1  = rnormal()
...
generate df y   = 1.0 + 2.0*x1 + 3.0*x2 + e

let m_ols   = ols(y ~ x1+x2+x3+x4+x5+x6+x7+x8+x9+x10, df)
let m_lasso = lasso(y ~ x1+x2+x3+x4+x5+x6+x7+x8+x9+x10, df, alpha=0.1)
let m_ridge = ridge_reg(y ~ x1+x2+x3+x4+x5+x6+x7+x8+x9+x10, df, alpha=0.5)
let m_enet  = enet(y ~ x1+x2+x3+x4+x5+x6+x7+x8+x9+x10, df, alpha=0.5)
\end{lstlisting}

\subsection*{OLS}

\begin{lstlisting}[language={}, basicstyle=\ttfamily\small,
  frame=none, numbers=none, backgroundcolor=\color{codbg},
  xleftmargin=0pt, xrightmargin=0pt, breaklines=false]
OLS: R²=0.801  (bom ajuste na amostra, mas x3..x10 poluem com coef não-nulos)
const=0.68  x1=1.93***  x2=2.88***  x3=0.17  x4=0.07  ...  x10=-0.11
\end{lstlisting}

\subsection*{LASSO}

\begin{lstlisting}[language={}, basicstyle=\ttfamily\small,
  frame=none, numbers=none, backgroundcolor=\color{codbg},
  xleftmargin=0pt, xrightmargin=0pt, breaklines=false]
LASSO α=0.1: R²=0.772  vars ativas: 2
  x1=1.552  x2=2.506  x3..x10 = 0.000  (zerando corretamente)
\end{lstlisting}

Com $\alpha = 0{,}1$ o LASSO identifica os dois preditores verdadeiros e
zera os oito irrelevantes.

\subsection*{Ridge}

\begin{lstlisting}[language={}, basicstyle=\ttfamily\small,
  frame=none, numbers=none, backgroundcolor=\color{codbg},
  xleftmargin=0pt, xrightmargin=0pt, breaklines=false]
Ridge α=0.5: R²=0.801
  x1=1.878  x2=2.793  x3=0.172  x4=0.081  ...
  (encolhe mas não zera nenhum)
\end{lstlisting}

\subsection*{Elastic Net}

\begin{lstlisting}[language={}, basicstyle=\ttfamily\small,
  frame=none, numbers=none, backgroundcolor=\color{codbg},
  xleftmargin=0pt, xrightmargin=0pt, breaklines=false]
ElasticNet α=0.5 l1_ratio=0.5: R²=0.589  vars ativas: 2
  x1=0.789  x2=1.550  x3..x10 = 0.000
\end{lstlisting}

Elastic Net também seleciona $x_1$ e $x_2$, mas com maior penalização total
(combinando $\ell_1$ e $\ell_2$) os coeficientes ficam mais encolhidos.

% ─────────────────────────────────────────────────────────────────────────────
\section{Sintaxe}

\begin{lstlisting}
lasso(y ~ x1 + x2 + x3, df, alpha=0.1)      // alpha = lambda de penalização
ridge_reg(y ~ x1 + x2 + x3, df, alpha=0.5)
enet(y ~ x1 + x2 + x3, df, alpha=0.5)       // alpha = mistura L1/(L1+L2)
\end{lstlisting}

\begin{center}
\begin{tabular}{lll}
\toprule
\textbf{Parâmetro} & \textbf{LASSO} & \textbf{Ridge / Enet} \\
\midrule
\texttt{alpha} & penalização $\lambda$ & penalização $\lambda$ \\
\texttt{alpha=0} & OLS & Ridge puro (Enet) \\
\texttt{alpha=1} & L1 máximo & L1 puro (Enet) \\
\bottomrule
\end{tabular}
\end{center}

\begin{nota}
  O valor ótimo de $\alpha$ é tipicamente escolhido por validação cruzada
  ($k$-fold). Em amostras pequenas, LASSO com $\alpha$ via \emph{BIC-LASSO}
  tende a ser mais conservador do que via CV.
\end{nota}
